

\section{Introduction\label{sec:introduction}}

This article compares {\it dynamic discrete choice\/} (DDC), a subfield of structural econometrics (SE),
to {\it inverse reinforcement learning\/} (IRL), a subfield of machine learning (ML). 
The goal of both is to infer preferences of individuals from observations of their choices over time 
assuming they are {\it rational,\/} i.e. that their choices over time maximize an expected reward (or ``utility'') function based on 
beliefs of how their actions affect uncertain current and future outcomes and payoffs. 
Individual  choices are affected by stochastic shocks that make their behavior appear probabilistic
from the standpoint of outside observers. DDC and IRL share a common framework, Markovian decision processes (MDP)
that can be viewed as dynamic versions of {\it random utility models\/} in economics, allowing
inference of preferences by methods such as
maximum likelihood, see \citet{mcfadden:1973} and \citet{rust:1987}. DDC and IRL differ primarily in the methods used to
solve the MDP, conduct inference, as well as the overall rationale for inferring preferences.

SE focuses on  testing and evaluating economic theories and making {\it counterfactual predictions\/}
of the effect of changes in the environment and
economic policies on economic behavior, outcomes, and welfare.\footnote{\footnotesize In economics ``policy'' refers to government interventions (tax rates, regulations, etc.), whereas
in RL/MDP it refers to a {\it decision rule,\/} i.e. a mapping of states into a probability distribution over actions.} 
SE differs from {\it reduced-form econometrics\/} that focus on summarizing behavior without attempting to infer 
preferences and thus without requiring the rationality assumption, analogous to the corresponding distinction
 between IRL and {\it imitation learning\/} (IL) in the ML literature.\footnote{\footnotesize
Other names for this include {\it behavioral cloning,\/} {\it learning from demonstration\/} and {\it apprenticeship learning.\/}
These are types of {\it supervised learning\/} in the ML literature, where an 
individual's decision rule can be learned from sufficient data on pairs of actions and states, $(a,s)$, without
any assumption that actions taken are ``optimal'',  see \citet{hussein:2018}.}
Reduced-form models can be used to make
some types of counterfactual predictions when there is sufficient
historical variation in policies to allow them to be treated as covariates in the model.
By uncovering preferences, SE models allow us to assess the welfare impacts of policy changes that are not
possible using reduced-form models and enable 
counterfactual predictions of policy changes that have no close historical antecedent.

For example, Denmark has long had one of the world's highest tax rates on new cars, but with little historical variation 
that could be used by reduced form approaches to predict how reductions in this tax would affect the Danish economy.
\citet{GIMRS:2022} structurally estimated an equilibrium model of ownership and trading of cars to recover consumer preferences
under the high tax {\it status quo\/} regime, and used it to compute counterfactual equilibria under a variety of alternative tax rates.
They showed that by reducing the tax on new cars and increasing the gas tax, Denmark can: 1) improve the welfare of most of its citizens,
2) increase total tax revenue from cars, and 3) reduce total $\mbox{CO}_2$ pollution.
Credible structural models are a faster, cheaper way of predicting the effects of policy changes by 
``crash testing'' them with computer simulations before actually implementing them in reality.

IRL, introduced by \citet{ng:2000}, builds on {\it reinforcement learning\/} (RL) \citep{BartoSutton:1998}. Its goal is to develop intelligent agents that can perform complex robotic and intellectual tasks without
the need for explicit programming or an explicit reward function to direct behavior. The success of ML 
is largely due to supervised learning algorithms (such as neural networks) to predict outcomes using labeled data. 
However this is data-intensive and does not naturally handle sequential decision making 
where actions have long-term consequences. RL overcomes this by training agents to maximize a reward through trial-and-error,
via repeated dynamic interactions with the environment. 
When the reward function is known, RL can  approximate optimal behavior without labeled examples. Chess is an example with a known reward (1 for a win, -1 for a loss and 0 for a draw). RL has been able to learn superhuman performance via {\it self-play,\/} i.e. repeatedly playing the algorithm against itself resulting in performance that exceeds the best human and computerized chess players, see e.g. \citet{Silver:2018}.

Unfortunately, it is often unclear what reward function to use to program tasks in AI and robotic applications using RL.
For example, \citet{christiano:2017} note that RL provides no guidance on the appropriate reward function to 
``train a robot to clean a table or scramble an egg.''  
IRL solves this by learning a reward function from limited data on human or expert demonstrations. Using the inferred reward function, 
these algorithms can be  trained more rapidly and cheaply via RL using computer simulations that result in
 intelligent behavior consistent with human preferences in situations outside of those in a limited demonstration
data set.
For example, Google Maps trip data \citep{Barnes:2024} revealed travelers' route preferences reflecting trade-offs between traffic, distance, hills, safety, and scenery. They massively scaled IRL to estimate a route preference function (360M parameters, 110M trips) enabling RL to generate recommendations achieving 16-24\% improvement in global route quality, and to the best of our knowledge, represents the largest published application of IRL in a real-world setting to date.

Section 2 reviews Dynamic Programming (DP) and Markovian Decision Processes (MDP), the common framework
used in the DDC and IRL literatures, and reinforcement learning (RL). Section 3 reviews the DDC literature
and section 4 the IRL literature. Section 5 concludes, comparing/contrasting
the literatures and the key challenges and limitations they both face.

%\begin{table}[h!]
%\centering
%\caption{List of abbreviations and acronyms}
%\begin{tabular}{|l|l|}
%\hline
%\textbf{Abbreviation} & \textbf{Full Name} \\
%\hline
%AI & Artificial Intelligence \\
%AIRL & Adversarial Inverse Reinforcement Learning\\
%AGI & Artificial General Intelligence \\
%AL & Apprenticeship Learning \\
%BC & Behavioral Cloning \\
%BIRL & Bayesian Inverse Reinforcement Learning\\
%CCP & Conditional Choice Probability \\
%CCS & Conditional Choice Simulator \\
%DDC & Dynamic Discrete Choice \\
%DM  & Decision Maker (Agent) \\
%DNN & Deep Neural Network \\
%DP  & Dynamic Programming \\
%DSE & Dynamic Structural Estimation \\
%DU & Discounted Utility \\
%EU   & Expected Utility\\
%IOC & Inverse Optimal Control \\
%IL & Imitation Learning  \\
%IRL & Inverse Reinforcement Learning\\
%MCMC & Markov Chain Monte Carlo \\
%MDP & Markovian Decision Process \\
%ML & Machine Learning \\
%MPE & Markov perfect Equilibrium \\
%MPEC & Mathematical Programming with Equilibrium Constraints\\
%NFXP & Nested Fixed Point \\
%%PG & Policy Gradient \\
%PI  & Policy Iteration \\
%%PPO & Proximal Policy Optimization \\
%RFE & Reduced Form Estimation \\
%RL & Reinforcement Learning  \\
%RLHF & Reinforcement Learning from Human Feedback \\
%RTDP & Real Time Dynamic Programming \\
%SA  & Successive Approximations \\
%TD  & Temporal Difference\\
%%TRPO & Trust Region Polcy Optimization \\
%%XAI, XRL & Explainable AI, RL \\
%\hline
%\end{tabular}
%\label{tab:abbrevs}
%\end{table}
